\chapter{Solving Linear Systems}
\label{ch:solving-linear-systems}

\begin{draftbox}
This chapter is scaffolded. It will absorb the current notes on solving
$Ax=b$, inverse matrices, solution sets, and the optional Hamming code project.
\end{draftbox}

\section{Guiding Question}

When does $Ax=b$ have a solution, and how many solutions can it have?

\section{Planned Core Ideas}

\begin{itemize}
  \item existence means $b\in\Img A$;
  \item uniqueness means $\Ker A=\{0\}$;
  \item if $x_0$ is one solution, then all solutions are
    $x_0+\Ker A$;
  \item matrix inverses are special cases of existence and uniqueness;
  \item inconsistent systems lead naturally to least squares.
\end{itemize}

\section{Planned Python Thread}

Classify small systems by rank and residual. Compare direct solution,
least-squares output, and parametrized solution sets.

\section*{Chapter Summary}
\addcontentsline{toc}{section}{Chapter Summary}

To be completed in the chapter-shaping pass.

\section*{Exercises}
\addcontentsline{toc}{section}{Exercises}

\subsection*{A. Theory}

\begin{exercise}
\exerciseid{8.A1}
Prove that if $\hat{x}$ minimizes $\norm{Ax-b}$, then the residual
$b-A\hat{x}$ is orthogonal to every column of $A$.
\end{exercise}

\begin{exercise}
\exerciseid{8.A2}
Prove that $\Ker(A^TA)=\Ker(A)$ for every real matrix $A$.
\end{exercise}

\subsection*{B. By Hand}

\begin{exercise}
\exerciseid{8.B1}
Project $b=(2,1,0)$ onto the line spanned by $a=(1,1,1)$ and compute the
residual. Check that the residual is orthogonal to $a$.
\end{exercise}

\begin{exercise}
\exerciseid{8.B2}
For
\[
  A=\begin{bmatrix}1&0\\1&1\\1&2\end{bmatrix},
  \qquad
  b=\begin{bmatrix}1\\2\\2\end{bmatrix},
\]
write down the normal equations and solve them by hand.
\end{exercise}

\subsection*{C. Python / NumPy}

\begin{exercise}
\exerciseid{8.C1}
For a small overdetermined matrix $A$, use \texttt{np.linalg.lstsq} and verify
numerically that \texttt{A.T @ (b - A @ x)} is close to zero.
\end{exercise}

\begin{exercise}
\exerciseid{8.C2}
For a full-column-rank matrix $A$, build
\[
  H=A(A^TA)^{-1}A^T
\]
using \texttt{np.linalg.solve}. Check numerically that $H^2=H$ and $H^T=H$.
\end{exercise}

\subsection*{D. Challenge}

\begin{exercise}
\exerciseid{8.D1}
Give an example of a matrix $A$ and a vector $b$ for which $Ax=b$ has no
solution. Then decompose $b$ as $b_{\col}+b_\perp$, where
$b_{\col}\in\col(A)$ and $b_\perp\in\Ker(A^T)$.
\end{exercise}

